\noindent\textbf{Elección del procedimiento.} En vuelo horizontal se impone $L=mg=qSC_L$, con $q=\rho V^2/2$. Para pérdida se usa $C_{L,\max}$; para el crucero se calcula otro $C_L$. Después la polar proporciona $C_D$, la definición de resistencia da $D$ y la potencia útil es $DV$.

\begin{enumerate}

\item Despejar la velocidad positiva de $mg=(\rho V_s^2/2)S C_{L,\max}$.

\[\fitmath{\mathrm{Vs}=\sqrt{\frac{2\cdot m\cdot g}{\rho\cdot S\cdot C_{L,\max}}}}\]

\[\fitmath{\mathrm{Vs}=\sqrt{\frac{2\cdot (1100)\cdot (9{,}81)}{(1{,}225)\cdot (16{,}2)\cdot (1{,}6)}}\approx 26{,}07116\,\mathrm{m/s}}\]

\item La presión dinámica debe corresponder a 55 m/s, no a la pérdida.

\[\fitmath{q=\frac{\rho\cdot {V}^{2}}{2}}\]

\[\fitmath{q=\frac{(1{,}225)\cdot {(55)}^{2}}{2}\approx 1852{,}812\,\mathrm{Pa}}\]

\item Despejar el coeficiente de sustentación de equilibrio.

\[\fitmath{C_L=\frac{m\cdot g}{q\cdot S}}\]

\[\fitmath{C_L=\frac{(1100)\cdot (9{,}81)}{(1852{,}8125)\cdot (16{,}2)}\approx 0{,}3595135\,}\]

\item Introducir ese coeficiente en la polar, sin confundir $C_L$ y $C_D$.

\[\fitmath{C_D=C_{D0}+K\cdot {C_L}^{2}}\]

\[\fitmath{C_D=(0{,}027)+(0{,}054)\cdot {(0{,}3595135024)}^{2}\approx 0{,}0339795\,}\]

\item Convertir el coeficiente adimensional en fuerza.

\[\fitmath{D=q\cdot S\cdot C_D}\]

\[\fitmath{D=(1852{,}8125)\cdot (16{,}2)\cdot (0{,}03397949775)\approx 1019{,}914\,\mathrm{N}}\]

\item La potencia aerodinámica necesaria vence esa resistencia a la velocidad de vuelo.

\[\fitmath{P=D\cdot V}\]

\[\fitmath{P=(1019{,}913739)\cdot (55)\approx 56095{,}26\,\mathrm{W}}\]

\end{enumerate}
